\documentclass[11pt,oneside]{article}

\usepackage[T1]{fontenc}
\usepackage{lmodern}
\usepackage{microtype}
\usepackage[margin=1in,headheight=14pt]{geometry}
\usepackage{amsmath,amssymb,amsthm,mathtools}
\usepackage{mathrsfs}
\usepackage{booktabs,longtable,tabularx,array}
\usepackage{enumitem}
\usepackage{xcolor}
\usepackage{xurl}
\usepackage{hyperref}
\usepackage{fancyhdr}
\usepackage{titlesec}
\usepackage{listings}
\usepackage{needspace}

\definecolor{proofgreen}{RGB}{17,105,72}
\definecolor{classicalblue}{RGB}{32,82,147}
\definecolor{targetviolet}{RGB}{112,54,145}
\definecolor{computeorange}{RGB}{174,87,0}
\definecolor{warningred}{RGB}{150,36,45}
\definecolor{codegray}{RGB}{248,248,248}
\definecolor{rulegray}{RGB}{105,110,118}

\hypersetup{
  colorlinks=true,
  linkcolor=classicalblue,
  citecolor=proofgreen,
  urlcolor=classicalblue,
  pdftitle={Beyond the Unified Report: Kummer-Symbol Separators and Output-Sensitive Newton Calculus},
  pdfauthor={Research continuation prepared for the ProveIt project},
  pdfsubject={New partial results and exact obstruction theorems for the Alaoglu--Erdos conjecture}
}
\urlstyle{same}

\pagestyle{fancy}
\fancyhf{}
\lhead{Alaoglu--Erd\H{o}s: new progress}
\rhead{Kummer separators and Newton calculus}
\cfoot{\thepage}
\renewcommand{\headrulewidth}{0.4pt}

\titleformat{\section}{\Large\bfseries}{\thesection}{0.75em}{}
\titleformat{\subsection}{\large\bfseries}{\thesubsection}{0.75em}{}
\titleformat{\subsubsection}{\normalsize\bfseries}{\thesubsubsection}{0.75em}{}

\setlength{\parindent}{0pt}
\setlength{\parskip}{0.55em}
\setlength{\emergencystretch}{3em}
\setlist[itemize]{leftmargin=2em,itemsep=0.25em,topsep=0.35em}
\setlist[enumerate]{leftmargin=2.25em,itemsep=0.25em,topsep=0.35em}

\newtheorem{theorem}{Theorem}[section]
\newtheorem{proposition}[theorem]{Proposition}
\newtheorem{lemma}[theorem]{Lemma}
\newtheorem{corollary}[theorem]{Corollary}
\newtheorem{conjecture}[theorem]{Conjecture}
\newtheorem{criterion}[theorem]{Criterion}
\newtheorem{question}[theorem]{Question}
\theoremstyle{definition}
\newtheorem{definition}[theorem]{Definition}
\newtheorem{experiment}[theorem]{Computational experiment}
\newtheorem{target}[theorem]{Research target}
\newtheorem{warning}[theorem]{Caution}
\theoremstyle{remark}
\newtheorem{remark}[theorem]{Remark}

\newcommand{\N}{\mathbb N}
\newcommand{\Nzero}{\mathbb N_0}
\newcommand{\Z}{\mathbb Z}
\newcommand{\Q}{\mathbb Q}
\newcommand{\R}{\mathbb R}
\newcommand{\F}{\mathbb F}
\newcommand{\Gm}{\mathbb G_m}
\newcommand{\Gal}{\operatorname{Gal}}
\newcommand{\ord}{\operatorname{ord}}
\newcommand{\lcm}{\operatorname{lcm}}
\newcommand{\rad}{\operatorname{rad}}
\newcommand{\supp}{\operatorname{supp}}
\newcommand{\Span}{\operatorname{span}}
\newcommand{\dens}{\operatorname{dens}}
\newcommand{\vpp}[2]{v_{#1}\!\left(#2\right)}
\newcommand{\dd}{\mathop{\mathrm d}}
\newcommand{\AEname}{Alaoglu--Erd\H{o}s}
\newcommand{\thetaAE}{\vartheta}
\newcommand{\statusproved}{\textcolor{proofgreen}{\textbf{[proved here]}}}
\newcommand{\statusclassical}{\textcolor{classicalblue}{\textbf{[classical input]}}}
\newcommand{\statuscompute}{\textcolor{computeorange}{\textbf{[exact computation]}}}
\newcommand{\statustarget}{\textcolor{targetviolet}{\textbf{[open target]}}}
\newcommand{\statuswarning}{\textcolor{warningred}{\textbf{[scope warning]}}}

\newenvironment{statusbox}[1]
{\begin{center}\begin{minipage}{0.94\textwidth}\hrule\vspace{0.45em}\textbf{Status.} #1\par\vspace{0.35em}}
{\vspace{0.35em}\hrule\end{minipage}\end{center}}

\newenvironment{resultbox}
{\begin{center}\begin{minipage}{0.94\textwidth}\hrule\vspace{0.45em}}
{\vspace{0.35em}\hrule\end{minipage}\end{center}}

\lstdefinestyle{pythonstyle}{
  language=Python,
  basicstyle=\ttfamily\scriptsize,
  backgroundcolor=\color{codegray},
  frame=single,
  rulecolor=\color{rulegray},
  numbers=left,
  numberstyle=\tiny\color{rulegray},
  numbersep=7pt,
  showstringspaces=false,
  breaklines=true,
  columns=fullflexible,
  keepspaces=true,
  tabsize=4
}

\begin{document}
\hypersetup{pageanchor=false}
\pagenumbering{gobble}

\begin{titlepage}
\centering
\vspace*{0.7cm}
{\Huge\bfseries Beyond the Unified Report\par}
\vspace{0.35cm}
{\LARGE\bfseries Kummer-Symbol Separators and\par}
\vspace{0.12cm}
{\LARGE\bfseries Output-Sensitive Newton Calculus\par}
\vspace{0.6cm}
{\Large A new progress report on the Alaoglu--Erd\H{o}s conjecture\par}
\vspace{0.8cm}
{\large Research continuation prepared for the ProveIt project\par}
\vfill
\begin{minipage}{0.9\textwidth}
\small
\textbf{Bottom line.}
No unconditional proof of
\[
  2^x,3^x\in\Z\quad\Longrightarrow\quad x\in\Z
\]
is claimed here.  The attached unified report already exhausts a remarkably broad family of
archimedean determinants, tropical and $p$-adic termwise divisors, sparse interpolation,
Sturmian dynamics, Loewner and operator formulations, canonical products, and standard
transcendence reductions.  This continuation therefore attacks a different interface.

The main new theorem attaches to every non-power pair $(M,A)$ a finite Kummer determinant
whose nonvanishing occurs on a positive-density Chebotarev set of rational primes.  Such
primes cannot divide a simultaneous approximation gap
$\gcd(M^q-2^p,A^q-3^p)$ whenever one fixed prime $r$ does not divide $q$.  Thus any successful
gcd amplification must avoid a positive-density Frobenius set on infinitely many convergent
denominators.  Independently, an iterated geometric divided-difference calculus factors every
mixed Newton functional on the $2$--$3$ grid into explicit shifted gaps
$M-2^j$ and $A-3^j$.  Pairing complementary Newton words makes the synchronized numerator
independent of the schedule, and a Kummer separator forbids simultaneous cancellation in
every fixed exponent residue class.  Chebotarev then gives a quantitative, though only
linear, lower bound for the surviving simultaneous denominator.

These statements are exact and prove useful new obstruction theorems.  They do not yet
connect the real logarithmic rank drop to enough finite-place divisibility to cross the
four-exponentials barrier.  The report ends with a precise separator-hitting criterion that
would close the conjecture and a Lean-oriented decomposition of the finite and global parts.
\end{minipage}
\vfill
{\large 22 August 2026\par}
\end{titlepage}

\begin{abstract}
Let $M=2^x$ and $A=3^x$ be positive integers.  The conjecture asks whether the real rank-one
condition
\[
  \log M\,\log3-\log A\,\log2=0
\]
forces $(M,A)=(2^n,3^n)$.  We give two new, complementary packages.

First, for an odd prime $r$ and a rational prime $\ell\equiv1\pmod r$, let
$\kappa_{\ell,r}:\F_\ell^*\to\F_r$ be an $r$th-power residue character and set
\[
 \Delta_{\ell,r}(M,A)=
 \kappa_{\ell,r}(2)\kappa_{\ell,r}(A)
 -\kappa_{\ell,r}(3)\kappa_{\ell,r}(M).
\]
We prove, using only finite linear algebra, Kummer theory, and Chebotarev, that unless
$(M,A)=(2^n,3^n)$ there is an odd $r$ for which
$\Delta_{\ell,r}(M,A)\ne0$ on a set of primes of positive natural density; an explicit
lower bound is $(r-2)/(r(r-1))$.  It follows that those primes divide no common gap
$\gcd(M^q-2^p,A^q-3^p)$ with $r\nmid q$.  This is a new support restriction on the exact gcd
whose exponential growth would be needed by rational-approximation approaches.

Second, for
$\mathcal N_{q,n}f(z)=[z,qz,\ldots,q^nz]f$ we prove the mixed word factorization
\[
 (\mathcal N_{q_h,n_h}\cdots\mathcal N_{q_1,n_1}t^x)(1)
 =\prod_{i=1}^h
  \frac{\prod_{j=0}^{n_i-1}(q_i^{x-S_{i-1}}-q_i^j)}
       {\prod_{j=0}^{n_i-1}(q_i^{n_i}-q_i^j)},
 \qquad S_i=n_1+\cdots+n_i.
\]
For $q_i\in\{2,3\}$ every numerator factors into explicit integer gaps.  Complementary words
produce the same synchronized product
$\prod_{j<N}(M-2^j)(A-3^j)$, independent of scheduling.  A residue-class version combines
with the Kummer theorem to show that positive-density separator primes cannot be cancelled in
both the dyadic and triadic Newton numerators.  The resulting Chebotarev mass is linear in the
Newton order, which is rigorous new progress but remains below the quadratic arithmetic scale
needed for closure.
\end{abstract}

\clearpage
\hypersetup{pageanchor=true}
\pagenumbering{arabic}
\tableofcontents
\clearpage

\section{Executive answer and claim ledger}
\label{sec:executive}

\subsection{What has and has not been proved}

The conjecture remains open.  This report does not promote a numerical pattern, a heuristic
independence assumption, or a conditional transcendence statement to a proof.  Its strongest
unconditional conclusions are the following.

\begin{resultbox}
\textbf{Positive-density Kummer separators.}
If $M,A\in\N$ are not simultaneously $(2^n,3^n)$, then for some odd prime $r$ a positive-density
set of rational primes $\ell\equiv1\pmod r$ satisfies
\[
 \Delta_{\ell,r}(M,A)\ne0.
\]
For every such $\ell$, the reduction of $(M,A)$ is not in the cyclic subgroup generated by
$(2,3)$ in $(\F_\ell^*)^2$.  The density is at least
$(r-2)/(r(r-1))$.  \statusproved{} with standard Kummer and Chebotarev inputs
\statusclassical.
\end{resultbox}

\begin{resultbox}
\textbf{A fixed-modulus exclusion for simultaneous gaps.}
For the same $r$ and the same Chebotarev set, if $r\nmid q$ then no separator prime divides
\[
  \gcd(M^q-2^p,A^q-3^p).
\]
More generally, a separator prime cannot divide both $M^q-2^p$ and $A^q-3^{p'}$ when
$r\nmid q$ and $p\equiv p'\pmod r$.  \statusproved
\end{resultbox}

\begin{resultbox}
\textbf{Exact mixed Newton factorization.}
Every ordered composition of geometric divided-difference operators on $t^x$ factors into
one-dimensional shifted-gap products.  For a counterexample pair, all factors are explicit
integers $M-2^j$ or $A-3^j$ divided by explicit geometric Vandermonde factors.  Pairing a word
with its base-complement makes the synchronized numerator independent of the word.
\statusproved
\end{resultbox}

\begin{resultbox}
\textbf{Quantitative simultaneous denominator survival.}
For the residue-class Newton factors with ratios $2^r$ and $3^r$, separator primes cannot
cancel both numerators.  Since every separator prime $\ell\le n+1$ divides both geometric
Vandermonde denominators of order $n$, the least common multiple of the two reduced
denominators has logarithm at least
\[
  (\delta_r+o(1))n,
\]
where $\delta_r>0$ is the Chebotarev density of the separator set.  \statusproved{} using the
Chebotarev prime-number theorem \statusclassical.
\end{resultbox}

The last lower bound is deliberately not oversold.  A linear surviving mass is too small to
settle the critical interpolation problem, whose natural denominator and height budgets are
quadratic or cubic in the relevant order.  Its value is structural: it proves that the
output-specific cancellation sought at the end of the unified report cannot be total, and it
identifies a concrete Frobenius-distribution theorem whose strengthening would matter.

\subsection{Novelty boundary relative to the attached report}

The source report is approximately $114{,}000$ words long and already reconciles twenty-seven
independent continuations.  Repeating its main reductions would obscure rather than advance
the project.  The following table states the boundary used here.

\begin{center}
\small
\begin{tabularx}{0.97\textwidth}{>{\raggedright\arraybackslash}p{0.31\textwidth}X}
\toprule
\textbf{Inherited and not re-proved} & \textbf{What is new in this continuation} \\
\midrule
Four Exponentials reduction; Gelfond--Schneider; Six Exponentials interfaces &
A finite Kummer-symbol determinant with a positive-density nonvanishing theorem. \\
Exact conditional monoid and primitive output normalization &
Use only the notation $(M,A)$; no duplicate classification or counting theory. \\
Ordinary and mixed generalized Vandermonde determinants; tropical ceilings &
Nested geometric Newton operators, order-sensitive factorization, and complementary-word invariance. \\
Iwasawa logarithm matrices, Fermat quotients, structural-prime transversality &
Good-prime $r$th-power residue characters distributed by Chebotarev. \\
Continued-fraction transfer and its unresolved common-gcd requirement &
A positive-density forbidden support theorem for that exact common gcd. \\
Canonical-product/Stieltjes denominator target &
A residue-class Newton denominator theorem showing unavoidable simultaneous denominator survival. \\
Loewner, operator, Sturmian, carry-language, entropy, Pad\'e, and Mahler audits &
No repetition; only short cross-checks where they delimit the new route. \\
\bottomrule
\end{tabularx}
\end{center}

\subsection{External-status note}

Anthropic publicly listed Claude Fable 5 and later announced its redeployment in June and
July 2026~\cite{AnthropicFable,AnthropicRedeploy}.  A search performed on 22 August 2026 did
not locate a publicly indexed manuscript, Lean repository, or precise theorem statement for
a Fable proof of the \AEname{} conjecture.  The rival team's message is therefore treated as
private, unverified information rather than as mathematical input.  Nothing below depends on
it.

\begin{statusbox}{The current-events paragraph is provenance only.  Every mathematical claim
in the report is independently stated and proved or clearly labeled as a classical input,
computation, or target.}
\end{statusbox}

\section{Minimal inherited setup}
\label{sec:setup}

\subsection{The counterexample notation}

Assume, only for the purpose of deriving consequences, that
\begin{equation}
  2^x=M\in\N,
  \qquad
  3^x=A\in\N.
  \label{eq:counterexample-pair}
\end{equation}
If $x<0$, both powers lie strictly between $0$ and $1$, so $x\ge0$.  If $x=0$ the conclusion
is immediate.  We therefore take $x>0$.  Taking real logarithms gives
\begin{equation}
  x=\frac{\log M}{\log2}=\frac{\log A}{\log3},
  \qquad
  \log M\,\log3-\log A\,\log2=0.
  \label{eq:real-determinant}
\end{equation}
The attached report proves that a nonintegral $x$ here would be transcendental and develops
its exact additive and multiplicative solution monoids.  We need none of that structure in
the first local theorem: the Kummer separator is valid for every positive integer pair
$(M,A)$ that is not a common integral power pair.

It is convenient to write
\[
 P=(2,3),\qquad Q=(M,A)\in \Gm^2(\Q).
\]
The desired conclusion is $Q=P^n$ for some $n\in\Nzero$, where powers are coordinatewise.
The real relation~\eqref{eq:real-determinant} says that the two vectors
$(\log2,\log3)$ and $(\log M,\log A)$ are parallel in $\R^2$.  Our local theorem asks whether
the reductions of $P$ and $Q$ are parallel in a finite cyclic quotient.  They are not: unless
$Q=P^n$, a positive-density set of primes detects a finite-place rank two.

\subsection{Prime-exponent vectors}

Let $S$ be the finite set of rational primes dividing $6MA$.  For a positive $S$-unit $u$,
write
\[
 \nu(u)=(v_p(u))_{p\in S}\in\Z^S.
\]
Set
\[
 e=\nu(2),\quad f=\nu(3),\quad m=\nu(M),\quad a=\nu(A).
\]
The vectors $e$ and $f$ are distinct standard basis vectors.  Unique factorization identifies
positive $S$-units modulo $r$th powers with the vector space
\[
 V_r=(\F_r)^S
\]
for any prime $r$.  A linear functional $\lambda\in V_r^*$ is therefore a formal
$r$th-power residue character on the prime support.

The central finite polynomial is
\begin{equation}
 \mathscr D_{M,A,r}(\lambda)
 =\lambda(e)\lambda(a)-\lambda(f)\lambda(m)\in\F_r.
 \label{eq:formal-kummer-det}
\end{equation}
It is the determinant of the two residue-symbol rows
$(\lambda(e),\lambda(f))$ and $(\lambda(m),\lambda(a))$.

\section{A positive-density Kummer-symbol separator theorem}
\label{sec:kummer}

\subsection{The formal rank-one criterion}

We first isolate the elementary linear algebra that makes the global argument work.

\begin{lemma}[Quadratic rank-one detector]
\label{lem:formal-rank-one}
Let $K$ be a field, let $V$ be a finite-dimensional $K$-vector space, and let
$e,f\in V$ be linearly independent.  For $m,a\in V$, suppose
\begin{equation}
 \lambda(e)\lambda(a)=\lambda(f)\lambda(m)
 \qquad\text{for every }\lambda\in V^*.
 \label{eq:all-functional-zero}
\end{equation}
Then there is $c\in K$ such that
\[
  m=ce,\qquad a=cf.
\]
Conversely, these two equalities imply~\eqref{eq:all-functional-zero}.
\end{lemma}

\begin{proof}
Choose $\lambda$ with $\lambda(e)=1$ and $\lambda(f)=0$.  Its values on a complement of
$\Span(e,f)$ are arbitrary.  Equation~\eqref{eq:all-functional-zero} says
$\lambda(a)=0$ for every such choice, so $a\in\Span(f)$.  Interchanging $e$ and $f$ gives
$m\in\Span(e)$.  Write $a=cf$ and $m=de$.  A functional nonzero on both $e$ and $f$ then
gives $c=d$.  The converse is immediate.
\end{proof}

\begin{corollary}[Formal determinant characterizes common powers]
\label{cor:formal-common-power}
Over $\Q$, the polynomial
\[
 \lambda\longmapsto
 \lambda(e)\lambda(a)-\lambda(f)\lambda(m)
\]
is identically zero if and only if
\[
  M=2^n,\qquad A=3^n
\]
for one $n\in\Nzero$.
\end{corollary}

\begin{proof}
Lemma~\ref{lem:formal-rank-one} gives $m=ce$ and $a=cf$ for $c\in\Q$.  The $2$-coordinate of
$m$ is the integer $v_2(M)=c$, so $c=n\in\Nzero$; the other coordinates vanish.  The same
argument applies to $A$.  The reverse implication is direct.
\end{proof}

\begin{remark}
This criterion is an algebraic tensor statement, not an independence conjecture for numerical
products of logarithms.  The latter was already isolated in the unified report.  Here the
coefficients are prime-exponent vectors and the variable is a finite residue character;
unique factorization makes the rank-one classification unconditional.
\end{remark}

\subsection{Reduction modulo an odd prime}

If $(M,A)$ is not a common power pair, Corollary~\ref{cor:formal-common-power} says that the
quadratic polynomial is nonzero over $\Q$.  Choose an odd prime $r$ for which its coefficient
tensor remains nonzero modulo $r$.  There are infinitely many such $r$.

A nonzero polynomial of total degree two on $(\F_r)^S$ cannot vanish everywhere when $r>2$.
The elementary Schwartz--Zippel bound gives
\begin{equation}
 \#\{\lambda\in V_r^*: \mathscr D_{M,A,r}(\lambda)=0\}
 \le 2r^{|S|-1}.
 \label{eq:schwartz-zippel}
\end{equation}
Thus at least $r^{|S|-1}(r-2)$ formal characters have nonzero determinant.

\subsection{Actual residue characters}

Let $\ell\nmid6MAr$ be a rational prime with $\ell\equiv1\pmod r$.  Choose a primitive
$r$th root of unity $\omega\in\F_\ell^*$.  The homomorphism
\begin{equation}
 \kappa_{\ell,r}:\F_\ell^*\longrightarrow\F_r,
 \qquad
 u^{(\ell-1)/r}=\omega^{\kappa_{\ell,r}(u)},
 \label{eq:residue-character}
\end{equation}
is an $r$th-power residue character.  Replacing $\omega$ by another primitive root scales
$\kappa_{\ell,r}$ by a nonzero element of $\F_r$.  Consequently the vanishing or nonvanishing
of
\begin{equation}
 \Delta_{\ell,r}(M,A)=
 \kappa_{\ell,r}(2)\kappa_{\ell,r}(A)
 -\kappa_{\ell,r}(3)\kappa_{\ell,r}(M)
 \label{eq:actual-kummer-det}
\end{equation}
is independent of that choice.

\begin{lemma}[Local cyclicity forces determinant zero]
\label{lem:local-cyclic-zero}
If the reduction of $Q=(M,A)$ belongs to the cyclic subgroup generated by
$P=(2,3)$ in $(\F_\ell^*)^2$, then $\Delta_{\ell,r}(M,A)=0$.
\end{lemma}

\begin{proof}
There is an integer $k$ such that $M\equiv2^k$ and $A\equiv3^k\pmod\ell$.  Applying the
homomorphism~\eqref{eq:residue-character} gives
\[
 \kappa(M)=k\kappa(2),\qquad \kappa(A)=k\kappa(3)
\]
in $\F_r$, and the determinant vanishes.
\end{proof}

\subsection{Kummer realization and Chebotarev density}

We now show that the nonzero formal characters occur as Frobenius residue characters on a
positive-density set of primes.

\begin{theorem}[Positive-density Kummer separators]
\label{thm:kummer-separators}
Let $M,A\in\N$.  If $(M,A)\ne(2^n,3^n)$ for every $n\in\Nzero$, then there is an odd prime
$r$ and a set $\mathcal P_r(M,A)$ of rational primes with positive natural density such that
for every $\ell\in\mathcal P_r(M,A)$:
\begin{enumerate}
\item $\ell\equiv1\pmod r$ and $\ell\nmid6MAr$;
\item $\Delta_{\ell,r}(M,A)\ne0$;
\item $Q\bmod\ell\notin\langle P\bmod\ell\rangle\subset(\F_\ell^*)^2$.
\end{enumerate}
One may choose $r$ so that
\begin{equation}
 \dens\mathcal P_r(M,A)\ge\frac{r-2}{r(r-1)}.
 \label{eq:density-lower}
\end{equation}
\end{theorem}

\begin{proof}
Choose an odd $r$ for which the formal determinant polynomial
$\mathscr D_{M,A,r}$ is nonzero.  Put $s=|S|$ and
$K=\Q(\zeta_r)$.  Consider the Kummer extension
\begin{equation}
 L=K\bigl(p^{1/r}:p\in S\bigr).
 \label{eq:kummer-field}
\end{equation}
The classes of the rational primes in $S$ are independent in $K^*/K^{*r}$.  Indeed, if
$\prod_{p\in S}p^{c_p}$ is an $r$th power in $K$ with $0\le c_p<r$, valuation at a prime of
$K$ above any $p\ne r$ gives $c_p\equiv0\pmod r$.  At the unique prime above $r$, the
valuation of $r$ is $r-1$, which is also coprime to $r$, so the same conclusion holds.
Therefore
\begin{equation}
 \Gal(L/K)\cong(\F_r)^S,
 \qquad [L:\Q]=r^s(r-1).
 \label{eq:kummer-galois}
\end{equation}
More explicitly, an element $\sigma_\lambda\in\Gal(L/K)$ indexed by
$\lambda\in(\F_r)^S$ acts by
\[
 \sigma_\lambda(p^{1/r})=\zeta_r^{\lambda_p}p^{1/r}.
\]

Let
\[
 \Omega_r=\{\lambda\in(\F_r)^S:
              \mathscr D_{M,A,r}(\lambda)\ne0\}.
\]
Conjugation by $\Gal(K/\Q)\cong\F_r^*$ scales $\lambda$ by a nonzero scalar, and the
quadratic determinant scales by its square.  Hence $\Omega_r$ is a union of conjugacy classes
in $\Gal(L/\Q)$.  Chebotarev's density theorem gives a set of unramified rational primes with
Frobenius in this union and natural density
\begin{equation}
 \frac{|\Omega_r|}{r^s(r-1)}.
 \label{eq:exact-cheb-density}
\end{equation}
Every such Frobenius restricts trivially to $K$, so the underlying rational prime satisfies
$\ell\equiv1\pmod r$.  The Kummer action on the radicals $p^{1/r}$ is exactly the
$r$th-power residue character~\eqref{eq:residue-character}, up to multiplication by a
nonzero scalar.  Thus its determinant is nonzero.  Lemma~\ref{lem:local-cyclic-zero} gives
local noncyclicity.

Finally,~\eqref{eq:schwartz-zippel} yields
$|\Omega_r|\ge r^{s-1}(r-2)$, and substitution in~\eqref{eq:exact-cheb-density} gives
\eqref{eq:density-lower}.
\end{proof}

\begin{statusbox}{The finite rank-one criterion and density count are elementary paper proofs.
The realization of characters uses standard Kummer theory and the density statement uses
Chebotarev~\cite{Washington,Neukirch}.  The theorem is compatible with, but more specialized
and more quantitative than, the classical support-problem literature
\cite{CorralesSchoof,PeruccaSupport}.  No claim of literature priority is made for this
specialization; the contribution here is the explicit determinant, density count, and its
common-gap bridge in the \AEname{} setting.}
\end{statusbox}

\begin{corollary}[Local-global characterization of common power pairs]
\label{cor:local-global-power}
For $M,A\in\N$, the following are equivalent:
\begin{enumerate}
\item $(M,A)=(2^n,3^n)$ for some $n\in\Nzero$;
\item for every odd prime $r$ and all but finitely many $\ell\equiv1\pmod r$,
      $\Delta_{\ell,r}(M,A)=0$;
\item there is no odd $r$ for which $\Delta_{\ell,r}(M,A)\ne0$ on a positive-density set of
      rational primes.
\end{enumerate}
\end{corollary}

\begin{proof}
The forward implication follows from multiplicativity of every residue character.  The
contrapositive of Theorem~\ref{thm:kummer-separators} gives the reverse implications.
\end{proof}

\subsection{What the theorem says about a hypothetical counterexample}

Under~\eqref{eq:counterexample-pair}, the real logarithm matrix has rank one.  If $x$ is
nonintegral, $Q$ is not an integral power of $P$, so Theorem~\ref{thm:kummer-separators}
produces an odd $r$ and positive-density primes at which the Kummer-symbol matrix has rank
two.  This is a genuine adelic asymmetry:
\begin{equation}
 \det\begin{pmatrix}\log2&\log M\\\log3&\log A\end{pmatrix}=0,
 \qquad
 \det\begin{pmatrix}\kappa(2)&\kappa(M)\\\kappa(3)&\kappa(A)\end{pmatrix}\ne0
 \quad\text{for a positive-density prime set}.
 \label{eq:real-kummer-asymmetry}
\end{equation}
It differs from the Iwasawa-logarithm transversality already present in the unified report in
two ways.  The primes are good primes rather than the structural primes $2$ and $3$, and the
nonvanishing is distributed by one explicit finite Galois extension rather than being a
pointwise local observation.

\section{The common-gap bridge}
\label{sec:common-gap}

The unified report identifies a sharp rational-approximation bottleneck.  If $p/q$ is close
to $x$, then the two integer gaps
\begin{equation}
 E_2(p,q)=M^q-2^p,
 \qquad
 E_3(p,q)=A^q-3^p
 \label{eq:common-gaps}
\end{equation}
are simultaneously small relative to $M^q$ and $A^q$.  To turn that archimedean closeness
into a contradiction, one needs unexpectedly large common arithmetic divisibility.  The new
Kummer theorem gives a fixed positive-density set that such common divisibility is forbidden
to use for most denominators $q$.

\subsection{A residue-symbol calculation}

\begin{proposition}[Mixed-gap determinant identity]
\label{prop:mixed-gap-det}
Fix an odd prime $r$ and a prime $\ell\equiv1\pmod r$ not dividing $6MA$.  Suppose
\begin{equation}
 M^q\equiv2^p\pmod\ell,
 \qquad
 A^q\equiv3^{p'}\pmod\ell.
 \label{eq:mixed-gap-congruences}
\end{equation}
Then
\begin{equation}
 q\,\Delta_{\ell,r}(M,A)
 =(p'-p)\,\kappa_{\ell,r}(2)\kappa_{\ell,r}(3)
 \quad\text{in }\F_r.
 \label{eq:mixed-gap-det-identity}
\end{equation}
In particular, if $r\nmid q$, $p\equiv p'\pmod r$, and
$\Delta_{\ell,r}(M,A)\ne0$, the two congruences in~\eqref{eq:mixed-gap-congruences} cannot
both hold.
\end{proposition}

\begin{proof}
Applying $\kappa=\kappa_{\ell,r}$ to~\eqref{eq:mixed-gap-congruences} gives
\[
 q\kappa(M)=p\kappa(2),
 \qquad
 q\kappa(A)=p'\kappa(3).
\]
Therefore
\begin{align*}
 q\Delta_{\ell,r}(M,A)
 &=\kappa(2)q\kappa(A)-\kappa(3)q\kappa(M)\\
 &=p'\kappa(2)\kappa(3)-p\kappa(3)\kappa(2),
\end{align*}
which is~\eqref{eq:mixed-gap-det-identity}.
\end{proof}

\begin{theorem}[Fixed-modulus exclusion from every common gap]
\label{thm:gap-exclusion}
Let $(M,A)$ not be a common integral power pair, and choose $r$ and
$\mathcal P_r(M,A)$ as in Theorem~\ref{thm:kummer-separators}.  For all integers $p$ and all
positive integers $q$ with $r\nmid q$,
\begin{equation}
 \ell\nmid \gcd(M^q-2^p,A^q-3^p)
 \qquad(\ell\in\mathcal P_r(M,A)).
 \label{eq:gap-exclusion}
\end{equation}
Equivalently, every integer
\[
 G_{p,q}(M,A)=\gcd(M^q-2^p,A^q-3^p),\qquad r\nmid q,
\]
is supported on the complement of one fixed positive-density Chebotarev set.
\end{theorem}

\begin{proof}
If $\ell$ divided both gaps, Proposition~\ref{prop:mixed-gap-det} with $p'=p$ would give
$q\Delta_{\ell,r}(M,A)=0$.  Since $r\nmid q$ and the determinant is nonzero, this is
impossible.
\end{proof}

\begin{remark}[Why the modulus is $r$, not $\ell-1$]
A direct cyclic-group argument would require $q$ to be invertible modulo the order of
$P\bmod\ell$, or at least modulo $\ell-1$.  The Kummer determinant only records the quotient
modulo $r$ and therefore needs the much weaker condition $r\nmid q$.  This fixed-modulus form
is essential for continued fractions: a fixed prime cannot divide every convergent
denominator.
\end{remark}

\subsection{Continued fractions and an unavoidable subsequence}

Let $p_j/q_j$ be the convergents of a nonintegral real $x$.  Consecutive denominators are
coprime:
\[
 \gcd(q_j,q_{j+1})=1.
\]
For the detector prime $r$, at least one member of each consecutive pair is not divisible by
$r$.  Consequently Theorem~\ref{thm:gap-exclusion} applies to infinitely many convergents,
with gaps that are simultaneously archimedean-small in the sense developed in the unified
report.

\begin{corollary}[Chebotarev avoidance along convergents]
\label{cor:convergent-avoidance}
If a nonintegral counterexample exists, there is an odd prime $r$, a positive-density
Chebotarev set $\mathcal P_r$, and an infinite set of indices $J$ such that
\begin{equation}
 \ell\nmid
 \gcd(M^{q_j}-2^{p_j},A^{q_j}-3^{p_j})
 \quad\text{for every }j\in J\text{ and every }\ell\in\mathcal P_r.
 \label{eq:convergent-exclusion}
\end{equation}
One can choose $J$ to meet every pair $\{j,j+1\}$.
\end{corollary}

This does not bound the numerical size of the gcd: one large integer may be composed entirely
of primes from a density-$1-\delta$ set.  It does, however, rule out any proof strategy that
would force these gcds to sample all Frobenius classes indiscriminately.

\subsection{A closure criterion}

The following criterion packages exactly what is missing.  It is deliberately stated in a
form that can be attacked by analytic or algebraic number theory without repeating the
archimedean transfer estimates.

\begin{criterion}[Separator-hitting criterion]
\label{crit:separator-hitting}
Suppose the following assertion holds.

For every positive integer pair $(M,A)$ satisfying
$\log M/\log2=\log A/\log3$, every odd prime $r$, and every positive-density Chebotarev set
$\mathcal C$ in a Kummer extension generated by $r$th roots of the primes dividing $6MA$,
there are integers $p,q$ with $r\nmid q$ and a prime $\ell\in\mathcal C$ such that
\[
 \ell\mid M^q-2^p,
 \qquad
 \ell\mid A^q-3^p.
\]
It is enough to require $p/q$ to range over convergents to
$\log M/\log2$.

Then the \AEname{} conjecture holds.
\end{criterion}

\begin{proof}
If $(M,A)$ were a non-power pair, take the detector prime $r$ and separator Chebotarev set
from Theorem~\ref{thm:kummer-separators}.  The asserted simultaneous divisibility contradicts
Theorem~\ref{thm:gap-exclusion}.
\end{proof}

\begin{warning}
Criterion~\ref{crit:separator-hitting} is not presently established and is stronger than a
mere lower bound for the size of the common gap gcd.  It demands Frobenius-class coverage.
The theorem's value is that this is now an exact, finite-extension target.  A proposed proof
can be tested against it without invoking an unspecified ``random prime divisor'' heuristic.
\end{warning}

\subsection{Relation to the classical support problem}

The support problem asks, roughly, whether inclusion or order divisibility of reductions at
almost all primes forces a global dependence relation.  For products of multiplicative
groups this has strong affirmative forms~\cite{CorralesSchoof,PeruccaSupport}.  Our direction
is the quantitative contrapositive specialized to the two points $P=(2,3)$ and $Q=(M,A)$:
a global failure of common-power dependence creates a positive-density set of local
separators.  The explicit determinant~\eqref{eq:actual-kummer-det} is useful because it also
interacts directly with approximation gaps through
Proposition~\ref{prop:mixed-gap-det}; a black-box support theorem alone would not expose the
fixed congruence $r\nmid q$.

\section{An output-sensitive geometric Newton calculus}
\label{sec:newton}

The unified report ends by isolating output-specific cancellation in the rank-one table
$M^aA^b$ as one of the few surviving interfaces.  Ordinary univariate interpolation on the
nodes $2^a3^b$ produces large, usually irreducible determinant numerators.  A different
operator ordering makes the output dependence completely factorized.

\subsection{One geometric block}

For distinct nodes $z_0,\ldots,z_n$, write
$[z_0,\ldots,z_n]f$ for the divided difference of order $n$.  For $q>1$ define
\begin{equation}
 \mathcal N_{q,n}f(z)=[z,qz,\ldots,q^nz]f,
 \qquad n\ge0.
 \label{eq:geometric-newton-operator}
\end{equation}
For $n=0$ this is just $f(z)$.

Put
\begin{equation}
 D_{q,n}=\prod_{j=0}^{n-1}(q^n-q^j)
 =q^{\binom n2}\prod_{h=1}^n(q^h-1),
 \qquad D_{q,0}=1.
 \label{eq:geometric-vandermonde}
\end{equation}

\begin{theorem}[Geometric divided difference of a real monomial]
\label{thm:one-block}
For $z>0$, $q>1$, $n\in\Nzero$, and $y\in\R$,
\begin{equation}
 \mathcal N_{q,n}(t^y)(z)
 =z^{y-n}
  \frac{\prod_{j=0}^{n-1}(q^y-q^j)}{D_{q,n}}.
 \label{eq:one-block-factorization}
\end{equation}
\end{theorem}

\begin{proof}
Scaling all nodes by $z$ extracts $z^{y-n}$, so take $z=1$.  The barycentric formula expresses
$[1,q,\ldots,q^n]t^y$ as a polynomial of degree at most $n$ in $R=q^y$.  It vanishes for
$R=1,q,\ldots,q^{n-1}$ because the $n$th divided difference of each polynomial
$t^k$, $k<n$, is zero.  Its leading coefficient is the reciprocal of
$\prod_{j=0}^{n-1}(q^n-q^j)$.  This determines the polynomial and gives
\eqref{eq:one-block-factorization}.
\end{proof}

\begin{remark}
The one-block identity is the standard geometric or $q$-Newton factorization.  The new point
below is to compose blocks of different bases, track the exponent shifts exactly, and pair
complementary words.  The order dependence is essential; these operators do not commute.
\end{remark}

\subsection{Mixed words}

Let
\[
 W=((q_1,n_1),\ldots,(q_h,n_h)),
 \qquad q_i>1,
 \quad n_i\in\N,
\]
and define cumulative orders
\[
 S_0=0,
 \qquad
 S_i=n_1+\cdots+n_i.
\]
The operator associated with $W$ is
\[
 \mathcal N_W=
 \mathcal N_{q_h,n_h}\circ\cdots\circ\mathcal N_{q_1,n_1}.
\]
Thus the first listed block acts first.

\begin{theorem}[Mixed-word factorization]
\label{thm:mixed-word}
For every real $x$,
\begin{equation}
 (\mathcal N_W t^x)(1)
 =\prod_{i=1}^h
 \frac{\displaystyle\prod_{j=0}^{n_i-1}
       \bigl(q_i^{x-S_{i-1}}-q_i^j\bigr)}
      {D_{q_i,n_i}}.
 \label{eq:mixed-word-factorization}
\end{equation}
More generally,
\[
 (\mathcal N_W t^x)(z)
 =z^{x-S_h}(\mathcal N_W t^x)(1).
\]
\end{theorem}

\begin{proof}
Theorem~\ref{thm:one-block} says that $\mathcal N_{q,n}$ maps the monomial $t^y$ to a scalar
multiple of $t^{y-n}$.  After the first $i-1$ blocks, the current exponent is
$x-S_{i-1}$.  Multiplying the scalar factors from Theorem~\ref{thm:one-block} proves the
formula by induction.
\end{proof}

\subsection{Specialization to a hypothetical counterexample}

Suppose~\eqref{eq:counterexample-pair} and each $q_i$ is either $2$ or $3$.  Put
\[
 B_2=M,
 \qquad
 B_3=A.
\]
Then
\begin{equation}
 q_i^{x-S_{i-1}}=q_i^{-S_{i-1}}B_{q_i}.
\end{equation}
and Theorem~\ref{thm:mixed-word} becomes the completely rational factorization
\begin{equation}
 (\mathcal N_W t^x)(1)
 =\prod_{i=1}^h
 q_i^{-S_{i-1}n_i}
 \frac{\displaystyle\prod_{j=0}^{n_i-1}
       \bigl(B_{q_i}-q_i^{S_{i-1}+j}\bigr)}
      {D_{q_i,n_i}}.
 \label{eq:counterexample-word-factorization}
\end{equation}
Every numerator factor is an explicit integer gap.  The exponent intervals
\[
 [S_{i-1},S_i)\cap\Z
\]
partition $\{0,1,\ldots,S_h-1\}$; the word labels each interval dyadic or triadic.

For two blocks one obtains, for example,
\begin{align}
 (\mathcal N_{3,m}\mathcal N_{2,n}t^x)(1)
 &=3^{-nm}
   \frac{\prod_{i=0}^{n-1}(M-2^i)}{D_{2,n}}
   \frac{\prod_{j=0}^{m-1}(A-3^{n+j})}{D_{3,m}},
 \label{eq:two-block-23}\\
 (\mathcal N_{2,n}\mathcal N_{3,m}t^x)(1)
 &=2^{-nm}
   \frac{\prod_{j=0}^{m-1}(A-3^j)}{D_{3,m}}
   \frac{\prod_{i=0}^{n-1}(M-2^{m+i})}{D_{2,n}}.
 \label{eq:two-block-32}
\end{align}
The shifted intervals make the noncommutativity transparent.

\subsection{Complementary-word invariance}

When every $q_i\in\{2,3\}$, let $\overline W$ have the same block lengths but replace every
$2$ by $3$ and every $3$ by $2$.

\begin{theorem}[Complementary-word product]
\label{thm:complementary-word}
Let $N=S_h$.  Under~\eqref{eq:counterexample-pair},
\begin{equation}
 (\mathcal N_Wt^x)(1)(\mathcal N_{\overline W}t^x)(1)
 =
 \frac{\displaystyle\prod_{e=0}^{N-1}(M-2^e)(A-3^e)}
 {\displaystyle
  6^{\sum_{i=1}^h S_{i-1}n_i}
  \prod_{i=1}^h D_{2,n_i}D_{3,n_i}}.
 \label{eq:complementary-product}
\end{equation}
The right-hand side depends on the ordered partition $(n_1,\ldots,n_h)$ but not on which
blocks of $W$ were labeled $2$ or $3$.
\end{theorem}

\begin{proof}
For a fixed interval $[S_{i-1},S_i)$, the two complementary factors in
\eqref{eq:counterexample-word-factorization} contribute every dyadic gap and every triadic
gap in that interval.  Their shift powers multiply to
$2^{-S_{i-1}n_i}3^{-S_{i-1}n_i}=6^{-S_{i-1}n_i}$, and their geometric denominators multiply
to $D_{2,n_i}D_{3,n_i}$.  Multiplication over the partition proves the formula.
\end{proof}

For unit blocks $n_i=1$, the formula simplifies to
\begin{equation}
 (\mathcal N_Wt^x)(1)(\mathcal N_{\overline W}t^x)(1)
 =
 \frac{\displaystyle\prod_{e=0}^{N-1}(M-2^e)(A-3^e)}
 {2^N6^{N(N-1)/2}}.
 \label{eq:unit-complement-product}
\end{equation}
The exact schedule independence is useful experimentally: dynamic programming over words can
optimize one functional without changing the paired synchronized numerator.

\subsection{Structural-prime valuations}

The factorization also gives exact $2$-adic and $3$-adic bookkeeping, rather than a termwise
minimum estimate.

\begin{proposition}[Exact structural valuation of a shifted block]
\label{prop:structural-valuation}
Let $q$ be prime, $B\in\N$, $a_0=v_q(B)$, and write $B=q^{a_0}u$ with $q\nmid u$.
For $S,n\in\Nzero$ with $n>0$ and $B\ne q^{S+j}$ for $0\le j<n$, put
\[
 N_{q,n,S}(B)=\prod_{j=0}^{n-1}(B-q^{S+j}).
\]
Then
\begin{equation}
 v_q\bigl(N_{q,n,S}(B)\bigr)
 =\sum_{j=0}^{n-1}\min(a_0,S+j)
 +\mathbf 1_{\{S\le a_0<S+n\}}v_q(u-1).
 \label{eq:exact-numerator-qvaluation}
\end{equation}
Consequently
\begin{equation}
 v_q\!\left(
 q^{-Sn}\frac{N_{q,n,S}(B)}{D_{q,n}}
 \right)
 =\sum_{j=0}^{n-1}\min(a_0,S+j)-Sn-\binom n2
 +\mathbf 1_{\{S\le a_0<S+n\}}v_q(u-1).
 \label{eq:exact-block-qvaluation}
\end{equation}
\end{proposition}

\begin{proof}
For $a_0\ne S+j$,
$v_q(B-q^{S+j})=\min(a_0,S+j)$.  At equality the valuation is
$a_0+v_q(u-1)$, accounting for the correction term.  Finally,
$v_q(D_{q,n})=\binom n2$ by~\eqref{eq:geometric-vandermonde}.
\end{proof}

Two clean regimes are worth recording:
\begin{align*}
 a_0<S:
 &\quad v_q(\text{block})=n(a_0-S)-\binom n2,\\
 a_0\ge S+n:
 &\quad v_q(\text{block})=0.
\end{align*}
Thus a block placed strictly beyond the $q$-adic depth of the output pays a predictable
quadratic denominator, while a block entirely below that depth is a $q$-adic unit.  This
turns word ordering into an exact finite optimization problem.

\subsection{Away-prime cancellation windows}

For a prime $\ell\nmid qB$,
\begin{equation}
 v_\ell(D_{q,n})=\sum_{h=1}^n v_\ell(q^h-1),
 \qquad
 v_\ell(N_{q,n,S}(B))
 =\sum_{e=S}^{S+n-1}v_\ell(B-q^e).
 \label{eq:away-valuations}
\end{equation}
At the first residue layer,
\begin{equation}
 \ell\mid N_{q,n,S}(B)
 \iff
 B\bmod\ell=q^e\bmod\ell
 \text{ for some }e\in[S,S+n)\cap\Z.
 \label{eq:discrete-log-window}
\end{equation}
Hence reduction of the Newton denominator is a discrete-log window problem.  The output
specificity absent from node-only Vandermonde bounds is now explicit: a prime can cancel only
if the actual output $B$ lands in a prescribed interval of the cyclic orbit of $q$.

\begin{statusbox}{Equations~\eqref{eq:counterexample-word-factorization},
\eqref{eq:complementary-product}, \eqref{eq:exact-block-qvaluation}, and
\eqref{eq:discrete-log-window} are the precise arithmetic interface exported by the Newton
calculus.  They do not by themselves make the functional small enough for a unit-threshold
contradiction.  Their role is to replace an opaque determinant numerator by explicit output
gaps whose cancellation can be tested prime by prime.}
\end{statusbox}

\section{Synthesis: residue-class Newton factors meet Kummer separators}
\label{sec:synthesis}

The Kummer determinant remembers exponents only modulo $r$.  The correct Newton factors to
pair with it therefore group gaps by one exponent residue class.

\subsection{Residue-synchronized gap products}

Fix the detector prime $r$ from Theorem~\ref{thm:kummer-separators}.  For
$c\in\Z/r\Z$ and $N\ge1$, define
\begin{align}
 U_{c,N}^{(r)}(M)
 &=\prod_{\substack{0\le e<N\\e\equiv c\pmod r}}(M-2^e),
 \label{eq:residue-dyadic-product}\\
 V_{c,N}^{(r)}(A)
 &=\prod_{\substack{0\le e<N\\e\equiv c\pmod r}}(A-3^e).
 \label{eq:residue-triadic-product}
\end{align}
Empty products are $1$.  Whenever $c$ occurs in an integer exponent below, we use the unique
representative $0\le c<r$.

\begin{theorem}[Residue-synchronized gcd sieve]
\label{thm:residue-gcd-sieve}
For every separator prime $\ell\in\mathcal P_r(M,A)$, every residue class $c\pmod r$, and
every $N$,
\begin{equation}
 \ell\nmid
 \gcd\bigl(U_{c,N}^{(r)}(M),V_{c,N}^{(r)}(A)\bigr).
 \label{eq:residue-gcd-exclusion}
\end{equation}
\end{theorem}

\begin{proof}
If $\ell$ divided both products, there would be exponents $e,f<N$ with
$e\equiv f\equiv c\pmod r$ and
\[
 M\equiv2^e\pmod\ell,
 \qquad
 A\equiv3^f\pmod\ell.
\]
Proposition~\ref{prop:mixed-gap-det} with $q=1$, $p=e$, and $p'=f$ would give
$\Delta_{\ell,r}(M,A)=0$, contrary to the definition of the separator set.
\end{proof}

The result is stronger than excluding the termwise gcds
$\gcd(M-2^e,A-3^e)$: a separator prime cannot divide a dyadic gap and a triadic gap at
\emph{different} exponents as long as the two exponents are congruent modulo $r$.

\subsection{Residue-class Newton numerators}

Let $Q>1$.  Applying Theorem~\ref{thm:one-block} with ratio $Q=q^r$ and exponent
$(x-c)/r$ gives
\begin{equation}
 [1,q^r,q^{2r},\ldots,q^{rn}]\,t^{(x-c)/r}
 =q^{-cn}
 \frac{\prod_{j=0}^{n-1}(B_q-q^{c+rj})}{D_{q^r,n}},
 \qquad B_2=M,\ B_3=A.
 \label{eq:residue-newton-factor}
\end{equation}
The values of $t^{(x-c)/r}$ at these nodes are
$(B_q/q^c)^j$, so after a transparent power of $q$ is cleared this is an arithmetic
functional of integers.  Its numerator is exactly one finite segment of the residue-class
products~\eqref{eq:residue-dyadic-product}--\eqref{eq:residue-triadic-product}.

For a fixed $c$ and $n$, put
\begin{align*}
 N_{2,c,n}&=\prod_{j=0}^{n-1}(M-2^{c+rj}),
 &F_{2,c,n}&=2^{-cn}\frac{N_{2,c,n}}{D_{2^r,n}},\\
 N_{3,c,n}&=\prod_{j=0}^{n-1}(A-3^{c+rj}),
 &F_{3,c,n}&=3^{-cn}\frac{N_{3,c,n}}{D_{3^r,n}}.
\end{align*}
Let $d_{2,c,n}$ and $d_{3,c,n}$ be their positive reduced denominators.

\begin{proposition}[A primewise simultaneous-denominator obstruction]
\label{prop:primewise-denominator}
Let $\ell\in\mathcal P_r(M,A)$.  If
$\ell\mid D_{2^r,n}$ and $\ell\mid D_{3^r,n}$, then
\begin{equation}
 v_\ell\!\left(\lcm(d_{2,c,n},d_{3,c,n})\right)
 \ge
 \min\bigl(v_\ell(D_{2^r,n}),v_\ell(D_{3^r,n})\bigr).
 \label{eq:lcm-denominator-lower-prime}
\end{equation}
\end{proposition}

\begin{proof}
Theorem~\ref{thm:residue-gcd-sieve} says that $\ell$ divides at most one of
$N_{2,c,n}$ and $N_{3,c,n}$.  Since $\ell\ne2,3$, the prefactors $2^{-cn}$ and $3^{-cn}$ do
not affect its valuation.  Therefore at least one reduced denominator retains the full
$\ell$-adic valuation of its corresponding geometric denominator.  This is at least the
minimum on the right side of~\eqref{eq:lcm-denominator-lower-prime}.
\end{proof}

\subsection{A quantitative Chebotarev mass}

Every prime $\ell\le n+1$ with $\ell\nmid6r$ divides both denominators in
Proposition~\ref{prop:primewise-denominator}.  Indeed,
$\ord_\ell(2^r)$ and $\ord_\ell(3^r)$ are at most $\ell-1\le n$, while
\[
 D_{Q,n}=Q^{\binom n2}\prod_{h=1}^n(Q^h-1).
\]
Thus the positive-density separator set creates a universal denominator residue.

\begin{theorem}[Linear simultaneous-denominator survival]
\label{thm:linear-denominator-mass}
Let $r$, $c$, and the reduced denominators $d_{2,c,n},d_{3,c,n}$ be as above, and let
$\delta_r=\dens\mathcal P_r(M,A)>0$.  Then
\begin{equation}
 \log\lcm(d_{2,c,n},d_{3,c,n})
 \ge (\delta_r+o(1))n
 \qquad(n\to\infty).
 \label{eq:linear-denominator-mass}
\end{equation}
In particular, with the construction in Theorem~\ref{thm:kummer-separators},
$\delta_r\ge(r-2)/(r(r-1))$.
\end{theorem}

\begin{proof}
For every $\ell\in\mathcal P_r(M,A)$ with $\ell\le n+1$, Proposition
\ref{prop:primewise-denominator} contributes at least one factor $\ell$ to the least common
multiple.  Hence
\[
 \log\lcm(d_{2,c,n},d_{3,c,n})
 \ge \sum_{\substack{\ell\le n+1\\\ell\in\mathcal P_r(M,A)}}\log\ell.
\]
The Chebotarev prime-number theorem evaluates the sum as
$(\delta_r+o(1))n$.
\end{proof}

\begin{remark}[Exact scale of the advance]
The full geometric denominators have
\[
 \log D_{q^r,n}=rn^2\log q+O(n),
\]
whereas Theorem~\ref{thm:linear-denominator-mass} forces only $\Omega(n)$ surviving mass.
Thus the theorem proves non-total cancellation but does not yet prevent a reduction by
$\exp(\Theta(n^2))$.  The missing quantitative problem is now concrete: estimate separator
primes whose orders of both $2^r$ and $3^r$ are at most $n$, not merely the automatic primes
$\ell\le n+1$.
\end{remark}

\begin{target}[Quadratic separator mass]
\label{target:quadratic-separator-mass}
For a detector Chebotarev set $\mathcal P_r(M,A)$, study
\begin{equation}
 \Sigma_r(n)=
 \sum_{\substack{\ell\in\mathcal P_r(M,A)\\
        \ord_\ell(2^r)\le n,\ \ord_\ell(3^r)\le n}}
 \min\bigl(v_\ell(D_{2^r,n}),v_\ell(D_{3^r,n})\bigr)\log\ell.
 \label{eq:separator-mass-sum}
\end{equation}
Theorem~\ref{thm:linear-denominator-mass} gives $\Sigma_r(n)\gg n$.  A lower bound
$\Sigma_r(n)\gg n^2$ would show that residue-class Newton denominators retain a positive
fraction of their natural arithmetic scale.  A sharper result coordinated with an
archimedean auxiliary function could close one of the surviving denominator-cancellation
interfaces.
\end{target}

This target is related to simultaneous multiplicative-order statistics for $2$ and $3$ in a
specified Frobenius set.  It is not a routine application of Artin's primitive-root
conjecture: the condition asks for two orders to be small rather than maximal, and the
primes are filtered by an independent Kummer determinant.  Nonetheless it is a much more
specific analytic problem than ``find extra cancellation in a determinant.''

\subsection{A finite local lemma suitable for formalization}

The global density theorem uses algebraic number theory, but the cancellation obstruction is
finite and can be isolated without Chebotarev.

\begin{lemma}[Finite residue-class obstruction]
\label{lem:finite-residue-class-obstruction}
Let $r$ be an odd prime, let $\ell\equiv1\pmod r$, and let
$\kappa:\F_\ell^*\to\F_r$ be a homomorphism.  Suppose
\[
 \kappa(2)\kappa(A)-\kappa(3)\kappa(M)\ne0.
\]
Then for any $c\in\F_r$ there do not exist integers $e,f\equiv c\pmod r$ with
$M\equiv2^e$ and $A\equiv3^f\pmod\ell$.
\end{lemma}

\begin{proof}
The two congruences imply
$\kappa(M)=e\kappa(2)$ and $\kappa(A)=f\kappa(3)$.  Their determinant is
$(f-e)\kappa(2)\kappa(3)=0$ because $e\equiv f\pmod r$.
\end{proof}

This lemma, Proposition~\ref{prop:mixed-gap-det}, and the mixed Newton factorization are the
parts most immediately amenable to machine checking.

\section{Scale audit: why the new identities do not yet prove the conjecture}
\label{sec:scale-audit}

A progress report on an open transcendence problem is useful only if it records the exact
point where a promising identity stops.  This section gives that audit.

\subsection{Archimedean size of one block}

For fixed real $y\notin\Nzero$ and fixed $q>1$, Theorem~\ref{thm:one-block} gives, as
$n\to\infty$,
\begin{equation}
 \log\left|\frac{\prod_{j=0}^{n-1}(q^y-q^j)}{D_{q,n}}\right|
 =-\frac12n^2\log q+O(n).
 \label{eq:block-arch-size}
\end{equation}
Indeed, all but $O_y(1)$ numerator factors have size $q^j(1+O(q^{y-j}))$, while
$\log D_{q,n}=n^2\log q+O(n)$.  The natural denominator of the rational specialization in
\eqref{eq:counterexample-word-factorization} has a quadratic logarithm of twice the leading
archimedean decay.  The bare one-block functional is therefore far above its reciprocal
denominator threshold.

For a word $W$, the leading decay is approximately
\begin{equation}
 -\frac12\sum_i n_i^2\log q_i,
 \label{eq:word-decay}
\end{equation}
whereas the explicit clearing denominator also contains the cross-shift tax
\begin{equation}
 \sum_i S_{i-1}n_i\log q_i.
 \label{eq:word-shift-tax}
\end{equation}
Splitting into many short blocks improves the sum of squares but increases the cross term;
using one long block does the reverse.  Without output-specific cancellation, neither
extreme crosses a unit threshold.  This recovers, in a new operator language, the
normalization barrier found abstractly in the unified report.

\subsection{What output cancellation would have to accomplish}

Let $d(W;M,A)$ be the reduced denominator of~\eqref{eq:counterexample-word-factorization}.
A successful Newton certificate would require an archimedean estimate of the form
\begin{equation}
 0<|\mathcal N_Wt^x(1)|<\frac1{d(W;M,A)}.
 \label{eq:newton-unit-threshold}
\end{equation}
The factorization reduces the arithmetic side of this inequality to valuations of
\[
 \prod_i\prod_{e=S_{i-1}}^{S_i-1}(B_{q_i}-q_i^e).
\]
At structural primes these are given exactly by Proposition
\ref{prop:structural-valuation}.  At away primes they are discrete-log windows
\eqref{eq:discrete-log-window}.  The Kummer theorem shows that the dyadic and triadic windows
cannot simultaneously cover a positive-density separator set when the exponents are aligned
modulo $r$.

Thus the new theory does not magically provide the cancellation; it converts the desired
cancellation into a falsifiable collection of local orbit-covering statements.

\subsection{Why linear separator mass is not enough}

The denominator lower bound~\eqref{eq:linear-denominator-mass} is a no-go statement for
perfect simultaneous reduction, but the critical budgets are quadratic.  If the reduced
denominator loses only $e^{cn}$ relative to an $e^{Cn^2}$ natural denominator, the leading
unit-threshold comparison is unchanged.  A proof through this route therefore needs one of
the following genuinely stronger inputs:
\begin{enumerate}
\item a quadratic lower bound for the separator mass~\eqref{eq:separator-mass-sum}, combined
      with a different auxiliary functional whose success requires excessive cancellation;
\item a quadratic \emph{upper} bound on allowed cancellation in the complement, using
      multiplicative-order or large-sieve information;
\item an archimedean construction with decay matching the already unavoidable linear
      denominator residue rather than the full quadratic scale;
\item a direct separator-hitting theorem for the continued-fraction common gaps, as in
      Criterion~\ref{crit:separator-hitting}.
\end{enumerate}
No one of these statements is supplied by generic Chebotarev alone.

\section{Additional directions explored and sharply delimited}
\label{sec:other-routes}

This section records several nonduplicate probes.  They are shorter because each hits a clear
ceiling before producing a theorem comparable to Sections~\ref{sec:kummer}--\ref{sec:synthesis}.

\subsection{Ordinary rectangular interpolation versus nested Newton factorization}

Take the four nodes $1,2,3,6$ with values $1,M,A,MA$.  The top ordinary univariate divided
difference is
\begin{equation}
 [1,2,3,6]f
 =\frac{MA-10A+15M-6}{60}.
 \label{eq:rect-2x2}
\end{equation}
The numerator is irreducible over $\Q[M,A]$.  Larger $2$--$3$ rectangles rapidly produce
higher-degree irreducible mixed numerators; see Section~\ref{sec:computations}.  In contrast,
the nested operator first takes a geometric difference along one scaling direction and then
along the other, yielding the completely factored expressions
\eqref{eq:two-block-23}--\eqref{eq:two-block-32}.

This comparison explains why the new operator ordering matters.  The ordinary top
coefficient treats the set $\{2^a3^b\}$ as unrelated one-dimensional nodes and entangles all
outputs.  The nested operator retains the multiplicative grid geometry and makes each
cancellation event local in the output.

\subsection{A rank-three prime-support probe}

Suppose, as a restricted model, that all prime factors of $MA$ lie in $\{2,3,p\}$.  Write
\[
 M=2^u3^vp^w,
 \qquad
 A=2^{u'}3^{v'}p^{w'}.
\]
The real rank-one relation becomes a quadratic relation among
$\log2,\log3,\log p$.  The two-prime-support theorem in the unified report handles the case
$w=w'=0$.  With one additional prime, elementary unique factorization no longer diagonalizes
the log-product relation.

The new Kummer determinant does diagonalize the corresponding finite problem.  The support
space has dimension three, and unless the exponent rows are a common integral multiple, a
quadratic form in three residue-character variables is nonzero.  Thus a positive-density
finite-place rank-two set exists even in the minimal unresolved support rank.  What remains
missing is a theorem that transfers the real equality to the vanishing of those residue
characters.  Standard Gelfond--Schneider and Six Exponentials inputs do not supply that
transfer, so this route merges back into the separator-hitting problem rather than producing
an independent closure.

\subsection{Cyclotomic symmetrization}

For a prime $q$, a tempting way to reduce a root-extraction error is the identity
\begin{equation}
 \Phi_{q^2}(z)=z^{q(q-1)}\frac{1-z^{-q^2}}{1-z^{-q}},
 \label{eq:cyclotomic-tail}
\end{equation}
whose normalized logarithmic correction is $O(z^{-q})$.  Applying it to radicals constructed
from $M,A,2,3$ gives excellent analytic cancellation.  However, clearing exponents and
integerizing the cyclotomic value introduces degree and height of order $q^2$, while the
gain in the logarithmic tail is only order $q$.  This is the same denominator/height tax
already diagnosed for high-root splittings in the unified report.  No new inequality emerges
without an output-specific common divisor.  The Kummer separator theorem is more informative:
it says that any proposed common divisor must avoid a specified positive-density set for a
fixed residue modulus.

\subsection{Subset-moment collisions}

Under a hypothetical counterexample, every number $t=2^aM^b$ and every shifted power
$t^{x+j}$ is an integer whenever $j\in\Nzero$.  Since the candidate set up to $e^T$ has
quadratic logarithmic cardinality, a subset-sum pigeonhole argument produces many signed
relations annihilating any fixed finite list of moments
\[
 t^0,t^1,\ldots,t^d,t^x,t^{x+1},\ldots,t^{x+d}.
\]
This is a genuine abundance phenomenon, but the relation support is of order $T^2$, while a
Chebyshev-system obstruction from $2d+2$ moments controls only relations with $O(d)$ sign
changes or support.  Forcing sparse collisions pays an entropy cost that removes the
pigeonhole surplus.  The route therefore does not presently beat the fixed-rank sum and
Chebyshev ceilings already recorded in the source report.

\subsection{Why the finite determinant cannot simply be ``taken to the limit''}

The real relation~\eqref{eq:real-determinant} and the finite determinant
\eqref{eq:actual-kummer-det} use different homomorphisms from the multiplicative group:
$\log|\cdot|$ is archimedean and continuous, while $\kappa_{\ell,r}$ is a torsion quotient
whose distribution varies with Frobenius.  There is no continuity principle taking one to
the other.  Any proof claiming that the real determinant forces
$\Delta_{\ell,r}=0$ for almost all $\ell$ would already establish a powerful new support
principle and must exhibit its arithmetic mechanism.  The common-gap bridge is one such
mechanism, but only at primes actually dividing both approximation gaps.

\begin{statusbox}{The explorations in this section are included to prevent repeated effort.
The principal new mathematical content remains the positive-density Kummer separator, the
fixed-modulus common-gap exclusion, and the mixed Newton/residue-class synthesis.}
\end{statusbox}

\section{Computational reconnaissance}
\label{sec:computations}

The computations in this section are exact symbolic arithmetic or elementary finite-field
enumeration.  They are diagnostics, not evidence for a theorem beyond the proved statements.
The core script used for the symbolic and exact-density tables is reproduced in Appendix~\ref{app:code}; the empirical prime loop is described there as well.

\subsection{Ordinary rectangular top coefficients}

For the rectangle
\[
 \mathcal R_{u,v}=\{2^a3^b:0\le a<u,\ 0\le b<v\}
\]
assign the value $M^aA^b$ at $2^a3^b$ and form the ordinary univariate top divided
difference.  Clearing content gives the following formulas.

\begin{align}
 \mathcal R_{2,2}:\quad
 &\frac{MA-10A+15M-6}{60},
 \label{eq:comp-22}\\[0.4em]
 \mathcal R_{3,2}:\quad
 &\frac{M^2A-66MA-880A+495M^2+594M-144}{47520},
 \label{eq:comp-32}\\[0.4em]
 \mathcal R_{2,3}:\quad
 &\frac{7MA^2-340A^2+1428MA-5712A+6885M-2268}{3084480}.
 \label{eq:comp-23}
\end{align}

Each displayed numerator is irreducible in $\Q[M,A]$.  For the $3\times3$ rectangle, the
primitive numerator is
\begin{align*}
 &M^2A^2-660MA^2-119680A^2+21420M^2A+565488MA\\
 &\qquad{}+2741760A-2044845M^2-1443420M+279936,
\end{align*}
with denominator $439723468800$; it is again irreducible over $\Q$.  The absence of a useful
factor is not mysterious: the ordinary
one-dimensional divided difference ignores the product-grid order.  The nested factorization
in Theorem~\ref{thm:mixed-word} restores it exactly.

\subsection{Exact finite Kummer densities in two examples}

For a fixed pair, the Kummer theorem predicts the density by counting the formal characters
$\lambda$.  Two small examples illustrate the calculation.

\begin{center}
\small
\begin{tabular}{@{}lcccc@{}}
\toprule
$(M,A)$ & $r$ & prime support $S$ &
$\Pr(\mathscr D(\lambda)\ne0\mid\ell\equiv1\bmod r)$ &
absolute density \\
\midrule
$(2,9)$ & $3$ & $\{2,3\}$ & $4/9$ & $2/9$ \\
$(5,13)$ & $3$ & $\{2,3,5,13\}$ & $16/27$ & $8/27$ \\
\bottomrule
\end{tabular}
\end{center}

For $(2,9)$ the determinant polynomial is $\lambda_2\lambda_3$, so exactly four of the nine
characters are separators.  For $(5,13)$ the four character values entering the $2\times2$
determinant are independent; the nonzero cases are the invertible $2\times2$ matrices over
$\F_3$, of which
\[
 |\operatorname{GL}_2(\F_3)|=(3^2-1)(3^2-3)=48.
\]
Dividing by $3^4$ gives $16/27$.  Multiplication by the cyclotomic splitting density $1/2$
gives the absolute densities.

Enumerating primes $\ell\le20{,}000$ with $\ell\equiv1\pmod3$ gives:

\begin{center}
\small
\begin{tabular}{@{}lrrrr@{}}
\toprule
$(M,A)$ & eligible primes & nonzero determinants & observed proportion & exact proportion \\
\midrule
$(2,9)$ & $1{,}124$ & $505$ & $0.44929$ & $0.44444\ldots$ \\
$(5,13)$ & $1{,}123$ & $672$ & $0.59840$ & $0.59259\ldots$ \\
\bottomrule
\end{tabular}
\end{center}

The experiment is only a check that the residue-symbol implementation and density formula
agree at a visible scale.

\subsection{Small-word checks}

An independent exact-rational checker enumerated every word of total order at most eight,
with every block size at most four and every base choice in $\{2,3\}$.  At three integral
exponents above the total order it compared the recursively evaluated nested divided
difference with~\eqref{eq:mixed-word-factorization} and also checked
complementary-word invariance~\eqref{eq:complementary-product}.  This gave $19{,}032$ exact
word checks.  A second loop performed $1{,}500$ exact valuation checks for
\eqref{eq:exact-block-qvaluation} at $q\in\{2,3,5\}$.

These checks are not used in the proofs.  They were useful for detecting the shift
$x-S_{i-1}$; omitting that shift is the most natural but incorrect way to compose the
operators.  A compact version of the checker is included after the main script in
Appendix~\ref{app:code}.

\section{Lean-oriented formalization plan}
\label{sec:lean}

The rival team's reported Lean activity makes it especially useful to separate the finite,
fully elementary core from the global Chebotarev input.  The following organization minimizes
the trusted interface.

\subsection{Module graph}

\begin{center}
\small
\begin{tabularx}{0.98\textwidth}{>{\ttfamily\raggedright\arraybackslash}p{0.32\textwidth}X}
\toprule
\textnormal{Module} & \textnormal{Contents} \\
\midrule
PrimeExponentVector & Finite support of a positive integer; exponent vector in
$S\to\Z$ and reduction to $S\to\operatorname{ZMod}r$. \\
FormalKummerDet & Definition of $\mathscr D_{M,A,r}$; Lemma~\ref{lem:formal-rank-one};
common-power classification; nonzero-polynomial witness over $\operatorname{ZMod}r$. \\
FiniteResidueCharacter & Multiplicative character
$\F_\ell^*\to\operatorname{ZMod}r$ and determinant invariance under rescaling. \\
LocalCyclicity & Lemma~\ref{lem:local-cyclic-zero}, Proposition~\ref{prop:mixed-gap-det},
and Theorem~\ref{thm:residue-gcd-sieve}. \\
GeometricDividedDifference & Barycentric divided differences, scaling, and
Theorem~\ref{thm:one-block}. \\
MixedNewtonWord & Cumulative orders, Theorems~\ref{thm:mixed-word} and
\ref{thm:complementary-word}. \\
BlockValuation & Proposition~\ref{prop:structural-valuation} and away-prime window lemmas. \\
KummerChebotarevInterface & Statement that formal characters are Frobenius characters with
the stated density.  Initially imported as one classical theorem; later expanded using
mathlib's algebraic-number-theory infrastructure where available. \\
SeparatorConsequences & Theorems~\ref{thm:kummer-separators},
\ref{thm:gap-exclusion}, and \ref{thm:linear-denominator-mass}. \\
\bottomrule
\end{tabularx}
\end{center}

\subsection{Finite statements first}

A schematic Lean statement for the key finite obstruction is:

\begin{lstlisting}[language={},basicstyle=\ttfamily\scriptsize,frame=single,backgroundcolor=\color{codegray}]
theorem mixed_gap_det
    {r ell : Nat} [Fact r.Prime] [Fact ell.Prime]
    (hr : r dvd ell - 1)
    (kappa : Units (ZMod ell) ->+ ZMod r)
    (hM : (M : ZMod ell) ^ q = (2 : ZMod ell) ^ p)
    (hA : (A : ZMod ell) ^ q = (3 : ZMod ell) ^ p') :
    q * kummerDet kappa M A =
      (p' - p) * kappa 2 * kappa 3 := by
  ...
\end{lstlisting}

The exact types should use additive notation for the target character, as in the
schematic declaration above.  The important point is dependency: this theorem requires
only finite cyclic groups and no analysis or algebraic number theory.

The mixed Newton theorem can be formalized by induction on a list of blocks after proving the
monomial eigenvector identity
\[
 \mathcal N_{q,n}(t^y)=C_{q,n}(y)t^{y-n}.
\]
For initial progress, one can state it for integer or rational symbolic exponents in a
Laurent-polynomial ring, then transfer the identity to real $x$ by direct evaluation.  The
counterexample specialization is purely algebraic after introducing hypotheses $2^x=M$ and
$3^x=A$.

\subsection{The global theorem as a narrow interface}

The Kummer/Chebotarev input can initially be isolated as follows.

\begin{lstlisting}[language={},basicstyle=\ttfamily\scriptsize,frame=single,backgroundcolor=\color{codegray}]
axiom exists_positive_density_kummer_character
    (hnotpower : not  exists  n : Nat, M = 2^n and A = 3^n) :
    exists  r : Nat, r.Prime and r != 2 and
      exists  P : Set Nat,
        HasNaturalDensity P delta and 0 < delta and
        forall  ell in P, exists  kappa, kummerDet kappa M A != 0
\end{lstlisting}

This is not suggested as the final trusted form; it is a staging boundary.  The paper proof of
Theorem~\ref{thm:kummer-separators} reduces it to four standard components:
\begin{enumerate}
\item independence of rational prime classes modulo $r$th powers in $\Q(\zeta_r)$;
\item the Kummer Galois-group computation;
\item conjugacy invariance of the nonzero determinant locus;
\item Chebotarev density for that union of conjugacy classes.
\end{enumerate}
Each component can be formalized independently.  Even before Chebotarev is available, all
finite consequences can be checked under an explicit Frobenius-character hypothesis.

\subsection{Suggested proof order}

The shortest machine-checked milestone is:
\begin{enumerate}
\item formal rank-one detector over an arbitrary field;
\item mixed-gap determinant identity over $\operatorname{ZMod}r$;
\item residue-synchronized gcd exclusion;
\item one-block and mixed-word Newton factorization;
\item complementary-word product;
\item structural valuation formula.
\end{enumerate}
This milestone captures every new algebraic identity in the report and leaves only density and
asymptotic prime summation outside the kernel.

\section{Prioritized next attacks}
\label{sec:next}

\subsection{Simultaneous-order Chebotarev mass}

The first priority is Target~\ref{target:quadratic-separator-mass}.  It admits both analytic
and computational attacks.  For each detector extension and each separator Frobenius class,
one can enumerate primes and record
\[
 \max(\ord_\ell(2^r),\ord_\ell(3^r)).
\]
The relevant cumulative weighted count is~\eqref{eq:separator-mass-sum}.  A useful theorem need
not resolve a full Artin-type conjecture; even a superlinear lower bound would quantify a new
barrier to simultaneous cancellation.

A possible analytic decomposition is by divisors:
\begin{equation}
 \mathbf 1_{\ord_\ell(2^r)\le n}
 \le \sum_{h\le n}\mathbf 1_{\ell\mid2^{rh}-1},
 \label{eq:order-divisor-decomp}
\end{equation}
with an analogous factor for $3$.  Their product leads to primes dividing resultants of
cyclotomic factors in two bases, filtered by a Kummer Frobenius class.  Generic upper bounds
are easy; the required lower bound is the hard direction.  The exact formulation nevertheless
permits targeted use of large-sieve-for-Frobenius methods rather than generic determinant
estimates.

\subsection{Common-gap Frobenius coverage}

The second priority is to estimate the Frobenius distribution of prime divisors of
\[
 G_j=\gcd(M^{q_j}-2^{p_j},A^{q_j}-3^{p_j}).
\]
Theorem~\ref{thm:gap-exclusion} says that a counterexample forces every $G_j$ with
$r\nmid q_j$ to avoid the separator classes.  An independent theorem saying that sufficiently
large $G_j$ must hit every class in the Kummer extension would combine directly with the
archimedean transfer machinery already present in the source report.

A promising technical form is not ``$G_j$ has a random prime divisor,'' but a lower bound for
its class-restricted radical:
\begin{equation}
 \log\rad_{\mathcal C}(G_j)
 =\sum_{\substack{\ell\mid G_j\\\operatorname{Frob}_\ell\in\mathcal C}}\log\ell.
 \label{eq:class-restricted-radical}
\end{equation}
Any positive lower bound for the separator class contradicts the exact exclusion.  This
formulation is compatible with sieve weights and avoids needing one specially selected prime.

\subsection{Word optimization with exact valuations}

The mixed Newton factorization turns auxiliary-function design into a finite scheduling
problem.  Given block lengths, each word assigns exponent intervals to bases $2$ and $3$.
The objective can combine:
\begin{enumerate}
\item exact structural valuations from Proposition~\ref{prop:structural-valuation};
\item away-prime cancellation from~\eqref{eq:discrete-log-window};
\item archimedean decay from~\eqref{eq:word-decay};
\item cross-shift tax from~\eqref{eq:word-shift-tax}.
\end{enumerate}
Dynamic programming can optimize over all words of total order hundreds or thousands without
expanding any determinant.  Complementary-word invariance supplies a consistency check and
allows one word to be optimized while the paired synchronized numerator stays fixed.

The computational goal is not to infer a proof from a favorable finite ratio.  It is to
identify a stable asymptotic scheduling law and then prove its valuation bounds symbolically.

\subsection{Search for a real-to-Kummer transfer}

The most direct route would prove that the real rank-one condition
\eqref{eq:real-determinant}, together with integrality of all four entries, forces
$\Delta_{\ell,r}=0$ on too many primes.  No known theorem does this.  Candidate mechanisms
include:
\begin{enumerate}
\item a new support theorem using archimedean proportionality as an additional place;
\item a reciprocity law for the Kummer determinant aggregated over primes;
\item an auxiliary algebraic number whose local symbols equal
      $\Delta_{\ell,r}(M,A)$ and whose archimedean height is forced to vanish;
\item a subspace-theorem construction in which nonzero Kummer determinant primes contribute a
      forbidden collection of local inequalities.
\end{enumerate}
The last two possibilities deserve particular attention because they could sum the
positive-density local obstruction into a global height contradiction.  The report has not
found the required auxiliary number.

\subsection{Formalization as a research instrument}

Formalizing the finite core is not merely presentational.  The identities involve three
moduli simultaneously: exponent shifts in $\Z$, residue exponents in $\F_r$, and reductions
in $\F_\ell$.  A proof assistant is well suited to preventing the common errors of inverting
$q$ modulo $\ell-1$ when only modulo $r$ is needed, confusing the two exponents $p,p'$, or
omitting the cumulative Newton shift $S_{i-1}$.  Once the interfaces are checked, automated
experiments can emit candidate lemmas in exactly the form needed by the formal development.

\section{Conclusion}
\label{sec:conclusion}

The old report's final diagnosis was correct: generic node geometry has already yielded most
of what it can yield.  The remaining arithmetic must see the actual outputs $M$ and $A$.
This continuation makes that statement concrete in two ways.

First, the Kummer-symbol determinant
\[
 \Delta_{\ell,r}(M,A)=
 \kappa(2)\kappa(A)-\kappa(3)\kappa(M)
\]
is an exact finite-place rank detector.  Unless $(M,A)$ is the desired common power pair, one
odd modulus $r$ produces nonzero determinant on a positive-density Chebotarev set.  Those
primes are excluded from every simultaneous approximation gap with $r\nmid q$.  Thus a
hypothetical counterexample must arrange its strongest common divisors inside the complement
of a fixed Frobenius set on an unavoidable subsequence of continued-fraction denominators.
This is a new, rigid global constraint.

Second, mixed geometric Newton operators factor exactly into shifted output gaps.  Their
complementary-word product synchronizes every dyadic and triadic gap while remaining
independent of the base schedule.  Grouping exponents modulo the detector prime couples this
factorization to the Kummer theorem: separator primes cannot cancel both residue-class
Newton numerators.  Chebotarev forces a linear amount of surviving simultaneous denominator.
The linear scale is not enough, but the previously qualitative phrase ``output-specific
cancellation'' has been replaced by the explicit mass sum~\eqref{eq:separator-mass-sum}.

The shortest credible path from here is therefore not another reformulation of Four
Exponentials.  It is one of two quantitative bridges:
\begin{enumerate}
\item prove that prime divisors of the continued-fraction common gaps hit every positive-density
      Kummer class; or
\item prove that separator primes contribute quadratic mass to the paired residue-class
      Newton denominators and combine this with a matching auxiliary-function estimate.
\end{enumerate}
Either theorem would interact immediately with exact identities already proved here.  Until
such a bridge is established, the mathematically honest status is substantial new local-global
progress, not a completed proof.

\appendix
\section{Details of the Kummer Galois computation}
\label{app:kummer-details}

For completeness, this appendix expands the two standard algebraic-number-theory facts used
in Theorem~\ref{thm:kummer-separators}.

\subsection{Independence of rational prime classes}

Let $r$ be odd and $K=\Q(\zeta_r)$.  Suppose
\[
 \prod_{p\in S}p^{c_p}=\alpha^r,
 \qquad \alpha\in K^*,
 \qquad 0\le c_p<r.
\]
For $p\ne r$, the extension $K/\Q$ is unramified at $p$.  At any prime ideal
$\mathfrak p$ above $p$,
\[
 v_{\mathfrak p}\!\left(\prod_{s\in S}s^{c_s}\right)=c_p.
\]
The right side $v_{\mathfrak p}(\alpha^r)$ is divisible by $r$, hence $c_p=0$.  At $p=r$,
there is one prime above $r$ and $v_{\mathfrak r}(r)=r-1$.  Thus
$c_r(r-1)$ is divisible by $r$; since $\gcd(r,r-1)=1$, again $c_r=0$.
This proves independence in $K^*/K^{*r}$ and hence
\[
 \Gal\bigl(K(p^{1/r}:p\in S)/K\bigr)\cong(\Z/r\Z)^S.
\]

\subsection{Conjugacy and density}

Let $N\cong(\F_r)^S$ be the Kummer subgroup and
$H\cong\F_r^*=\Gal(K/\Q)$.  Then
\[
 \Gal(L/\Q)\cong N\rtimes H,
\]
where $a\in H$ acts on $\lambda\in N$ by scalar multiplication $a\lambda$.  Since
$\mathscr D(a\lambda)=a^2\mathscr D(\lambda)$, the nonzero locus
$\Omega_r\subset N$ is stable under conjugacy in the full Galois group.  Chebotarev applied
to the union $\Omega_r$ gives density
\[
 \frac{|\Omega_r|}{|N\rtimes H|}
 =\frac{|\Omega_r|}{r^{|S|}(r-1)}.
\]
The Frobenius elements lie in $N$ and hence restrict trivially to $K$, which is equivalent to
$\ell\equiv1\pmod r$.  Choosing a prime of $K$ over $\ell$ chooses one representative
$\lambda$ in its scalar conjugacy orbit; determinant nonvanishing is independent of this
choice.

\section{Reproducibility code}
\label{app:code}

The following compact Python/SymPy script computes the rectangular divided differences and
the exact formal Kummer density.  The longer prime-enumeration loop used for the empirical
table differs only by adding \texttt{primerange}, \texttt{primitive\_root}, and
\texttt{discrete\_log}.

\begin{lstlisting}[style=pythonstyle]
from itertools import product
from collections import Counter
import sympy as sp

M, A = sp.symbols("M A")

def rectangular_top(u: int, v: int):
    points = []
    for a in range(u):
        for b in range(v):
            z = sp.Integer(2)**a * sp.Integer(3)**b
            y = M**a * A**b
            points.append((z, y))
    value = 0
    for i, (z, y) in enumerate(points):
        den = sp.prod(z - w for j, (w, _) in enumerate(points) if i != j)
        value += y / den
    return sp.together(value)

def exponent_vector(n: int, primes: list[int]) -> list[int]:
    fac = sp.factorint(n)
    return [fac.get(p, 0) for p in primes]

def formal_density(M0: int, A0: int, r: int):
    primes = sorted(sp.factorint(6*M0*A0))
    e = exponent_vector(2, primes)
    f = exponent_vector(3, primes)
    m = exponent_vector(M0, primes)
    a = exponent_vector(A0, primes)

    def dot(x, y):
        return sum(u*v for u, v in zip(x, y)) % r

    def det(lam):
        return (dot(lam, e)*dot(lam, a)
                - dot(lam, f)*dot(lam, m)) % r

    counts = Counter(det(lam)
                     for lam in product(range(r), repeat=len(primes)))
    nonzero = r**len(primes) - counts[0]
    conditional = sp.Rational(nonzero, r**len(primes))
    absolute = conditional / (r - 1)
    return primes, conditional, absolute

for shape in [(2, 2), (3, 2), (2, 3), (3, 3)]:
    print(shape, rectangular_top(*shape))

for pair in [(2, 9), (5, 13)]:
    print(pair, formal_density(*pair, r=3))
\end{lstlisting}

The exact word checker uses recursive barycentric evaluation rather than the factorization it
is testing.  Here is its compact core; the executed version additionally loops through all
compositions of total order at most eight and performs the valuation tests described in
Section~\ref{sec:computations}.

\begin{lstlisting}[style=pythonstyle]
from fractions import Fraction
from functools import lru_cache

def dd_geometric(f, q, n, z):
    xs = [z*q**j for j in range(n+1)]
    ans = Fraction(0)
    for i, xi in enumerate(xs):
        den = 1
        for j, xj in enumerate(xs):
            if i != j:
                den *= xi-xj
        ans += Fraction(f(xi), den)
    return ans

def direct_word(word, m):
    @lru_cache(None)
    def f(z):
        return Fraction(z**m)
    for q, n in word:
        old = f
        @lru_cache(None)
        def f(z, old=old, q=q, n=n):
            return dd_geometric(old, q, n, z)
    return f(1)

def D(q, n):
    return sp.prod(q**n-q**j for j in range(n))

def formula_word(word, m):
    ans = Fraction(1)
    shift = 0
    for q, n in word:
        num = sp.prod(q**(m-shift)-q**j for j in range(n))
        ans *= Fraction(int(num), int(D(q, n)))
        shift += n
    return ans

def check_complement(word, m):
    other = [(5-q, n) for q, n in word]
    lhs = formula_word(word, m)*formula_word(other, m)
    N = sum(n for _, n in word)
    shift = tax = 0
    den = 1
    for q, n in word:
        tax += shift*n
        den *= D(2, n)*D(3, n)
        shift += n
    M0, A0 = 2**m, 3**m
    num = sp.prod((M0-2**e)*(A0-3**e) for e in range(N))
    assert lhs == Fraction(int(num), int(6**tax*den))
\end{lstlisting}

\section{Glossary of the new interfaces}
\label{app:glossary}

\begin{longtable}{>{\raggedright\arraybackslash}p{0.27\textwidth}p{0.66\textwidth}}
\toprule
\textbf{Term} & \textbf{Meaning in this report} \\
\midrule
\endfirsthead
\toprule
\textbf{Term} & \textbf{Meaning in this report} \\
\midrule
\endhead
Kummer determinant & The finite determinant
$\kappa(2)\kappa(A)-\kappa(3)\kappa(M)$ in $\F_r$. \\
Separator prime & A good prime $\ell\equiv1\pmod r$ at which the Kummer determinant is
nonzero. \\
Detector modulus & An odd prime $r$ for which the formal determinant polynomial is nonzero
modulo $r$. \\
Common gap & The pair $M^q-2^p$, $A^q-3^p$, or their gcd. \\
Geometric Newton block & The divided-difference operator
$\mathcal N_{q,n}f(z)=[z,qz,\ldots,q^nz]f$. \\
Newton word & An ordered composition of geometric Newton blocks. \\
Complementary word & The word obtained by exchanging bases $2$ and $3$ while retaining block
lengths. \\
Residue-class numerator & A product of gaps $M-2^e$ or $A-3^e$ with $e$ in one class modulo
$r$. \\
Separator mass & The logarithmically weighted contribution of separator primes to geometric
Newton denominators; see~\eqref{eq:separator-mass-sum}. \\
Separator-hitting & A theorem forcing a common gap to acquire a prime divisor in a prescribed
positive-density Kummer Frobenius set. \\
\bottomrule
\end{longtable}

\begin{thebibliography}{99}

\bibitem{UnifiedReport}
ProveIt project,
\emph{The Alaoglu--Erd\H{o}s Integer-Exponent Conjecture: Unified Report},
unpublished working report, attached source version A, August 2026.

\bibitem{CorralesSchoof}
C.~Corrales-Rodrig\'a\~nez and R.~Schoof,
The support problem and its elliptic analogue,
\emph{Journal of Number Theory} \textbf{64} (1997), 276--290.

\bibitem{PeruccaSupport}
A.~Perucca,
Two variants of the support problem for products of abelian varieties and tori,
\emph{Journal of Number Theory} \textbf{129} (2009), 1883--1892;
preprint \url{https://arxiv.org/abs/0712.2815}.

\bibitem{Washington}
L.~C.~Washington,
\emph{Introduction to Cyclotomic Fields}, second edition,
Graduate Texts in Mathematics 83, Springer, 1997.

\bibitem{Neukirch}
J.~Neukirch,
\emph{Algebraic Number Theory},
Grundlehren der mathematischen Wissenschaften 322, Springer, 1999.

\bibitem{deBoor}
C.~de Boor,
Divided differences,
\emph{Surveys in Approximation Theory} \textbf{1} (2005), 46--69.

\bibitem{GasperRahman}
G.~Gasper and M.~Rahman,
\emph{Basic Hypergeometric Series}, second edition,
Encyclopedia of Mathematics and its Applications 96, Cambridge University Press, 2004.

\bibitem{BCZ}
Y.~Bugeaud, P.~Corvaja, and U.~Zannier,
An upper bound for the G.C.D. of $a^n-1$ and $b^n-1$,
\emph{Mathematische Zeitschrift} \textbf{243} (2003), 79--84.

\bibitem{AnthropicFable}
Anthropic,
\emph{Claude Fable 5}, product announcement and model page, 2026,
\url{https://www.anthropic.com/claude/fable}.

\bibitem{AnthropicRedeploy}
Anthropic,
\emph{Redeploying Fable 5}, 30 June 2026,
\url{https://www.anthropic.com/news/redeploying-fable-5}.

\end{thebibliography}

\end{document}
